\subsection{Examples of groupoids}

\begin{enumerate}
\def\labelenumi{\arabic{enumi})}
\item
  Let G be a group. Denote by \(\mathcal{G}\) the quiver whose set of vertices is reduced to a single element and whose set of arrows is G. Equipped with the composition law induced by that of G, \(\mathcal{G}\) is a groupoid, called the groupoid associated with the group G.
\item
  Let X be a set. Let G be the quiver \((\mathrm{X}, \mathrm{X} \times \mathrm{X}, \mathrm{pr}_{1}, \mathrm{pr}_{2})\) ; define a composition law on G by setting \((x, x') \cdot (x', x'') = (x, x'')\) , if \(x, x', x''\) are elements of X. This law is associative; for every \(x \in X\) , \((x, x)\) is the neutral element at x; the inverse of the arrow \((x, x')\) is the arrow \((x', x)\) . The groupoid thus defined is called the groupoid of pairs of the set X. If X is nonempty, it is simply transitive.
\item
  Let X and S be sets and let \(p \colon \mathrm{X} \to \mathrm{S}\) be a map. For \(a \in \mathrm{S}\) , write \(\mathrm{X}_a = p^{-1}(a)\) . For \(a\) and \(b\) in S, let \(\mathrm{G}_{a,b}\) be the set \(\mathcal{B}(\mathrm{X}_a; \mathrm{X}_b)\) of bijections from \(\mathrm{X}_a\) to \(\mathrm{X}_b\) and let G be the disjoint union of the sets \(\mathrm{G}_{a,b}\) . Let \(o\) and \(t\) be the maps from G to S such that \(o(f) = a\) and \(t(f) = b\) for every element \(f \in \mathrm{G}_{a,b}\) . The quadruple \((\mathrm{S}, \mathrm{G}, o, t)\) is a quiver. Equip it with the composition law defined by \(m(f, g) = g \circ f\) if \(f \in \mathcal{B}(\mathrm{X}_a; \mathrm{X}_b)\) and \(g \in \mathcal{B}(\mathrm{X}_b; \mathrm{X}_c)\) , where \(a, b, c\) are elements of S. This is a groupoid; it is denoted by \(\mathcal{B}(\mathrm{X}, p)\) and it is called the groupoid of permutations of X relative to \(p\) .
\item
  Let \((\mathrm{G}_{i})_{i\in\mathrm{I}}\) be a family of groupoids. Let G denote the product quiver of the family of underlying quivers; for every \(i\in I\) , let \(pr_{i}\colon G\to G_{i}\) be the canonical quiver morphism. There exists a unique composition law in G for which G is a groupoid and such that, for every \(i \in I\) , \(\mathrm{pr}_i\) is a morphism of groupoids. Let \(f = (f_i)_{i \in I}\) and \(g = (g_i)_{i \in I}\) be arrows of G; they are composable if and only if the arrows \(f_i\) and \(g_i\) are composable, for \(i \in I\) . One then has \(fg = (f_i g_i)_{i \in I}\) .
\end{enumerate}

The quiver G, equipped with this composition law, is called the product groupoid of the family \((\mathrm{G}_{i})_{i\in\mathrm{I}}\) . It satisfies the following universal property: for every groupoid \(G'\) and every family \((\varphi_{i})_{i\in\mathrm{I}}\) , where \(\varphi_{i}\) is a morphism of groupoids from \(G'\) to \(G_{i}\) , there exists a unique morphism of groupoids \(\varphi: G'\to G\) such that \(\varphi_{i}=\operatorname{pr}_{i}\circ\varphi\) for all \(i\in I\) .

\begin{enumerate}
\def\labelenumi{\arabic{enumi})}
\setcounter{enumi}{4}
\tightlist
\item
  Let \(((\mathrm{G}_{i})_{i\in\mathrm{I}},(\varphi_{i,j})_{i\prec j})\) be an inductive system of groupoids, indexed by a right-filtering preordered set \((\mathrm{I},\prec)\) , the \(\varphi_{i,j}\) being morphisms of groupoids. Let G be the quiver whose set of vertices is \(\varinjlim_{i}\operatorname{Som}(\mathbf{G}_{i})\) and whose set of arrows is \(\varinjlim_{i}\operatorname{Fl}(\mathbf{G}_{i})\) , the source and target maps being the maps \(\varinjlim_{i}o_{G_{i}}\) and \(\varinjlim_{i}t_{G_{i}}\) respectively. The composition laws of the \(G_{i}\) induce a composition law on G which makes it a groupoid (cf. A, I, p. 114). The canonical maps \(\operatorname{Som}(\mathbf{G}_{i})\to\operatorname{Som}(\mathbf{G})\) and \(\operatorname{Fl}(\mathbf{G}_{i})\to\operatorname{Fl}(\mathbf{G})\) define a morphism of groupoids \(\varphi_{i}\) from \(G_{i}\) to G. If i, j are elements of I such that \(i\prec j\) , we have \(\varphi_{j}\circ\varphi_{i,j}=\varphi_{i}\) . The groupoid G is called the inductive limit of the family of groupoids \(G_{i}\) and is denoted \(\varinjlim_{i}G_{i}\) . It satisfies the following universal property: for every groupoid H and every family \((\psi_{i})_{i\in\mathrm{I}}\) , where, for \(i\in I\) , \(\psi_{i}:G_{i}\to H\) is a morphism of groupoids, such that \(\psi_{j}\circ\varphi_{i,j}=\psi_{i}\) for every pair \((i,j)\) of elements of I such that \(i\prec j\) , there exists a unique morphism of groupoids \(\psi:G\to H\) such that \(\psi_{i}=\psi\circ\varphi_{i}\) .
\end{enumerate}

For the existence of a groupoid satisfying this universal property when one does not assume that the set I is filtering to the right, cf. II, p. 228, exerc. 3.
% semantic-fragments: legacy-input-seam
\relax{}
